\documentclass[UTF8]{ctexart}

\usepackage[a4paper, margin=2.2cm]{geometry}
\usepackage{booktabs}
\usepackage{graphicx}
\usepackage{float}
\usepackage{hyperref}
\usepackage{array}
\usepackage{tabularx}

\newcolumntype{Y}{>{\raggedright\arraybackslash}X}
\renewcommand{\arraystretch}{1.15}

\title{CIFAR-10 传统 CNN Baseline 实验记录}
\author{}
\date{}

\begin{document}
\maketitle

\section{模型结构}

本实验首先实现一个传统卷积神经网络作为 CIFAR-10 分类任务的 baseline。模型采用经典的
\texttt{Conv2d-ReLU-MaxPool2d-FullyConnected} 结构，不使用 Batch Normalization 和残差连接，
以便后续与更复杂的网络结构进行对比。模型输入为 $3\times32\times32$ 图像，输出为 10 类 logits。
可训练参数量为 666,538。

\begin{table}[H]
    \centering
    \caption{ClassicCNN 网络结构}
    \begin{tabularx}{\textwidth}{llY}
        \toprule
        模块 & 层设置 & 输出尺寸 \\
        \midrule
        Input & CIFAR-10 RGB image & $3\times32\times32$ \\
        Block 1 & Conv $3\rightarrow32$, $3\times3$; ReLU; Conv $32\rightarrow32$, $3\times3$; ReLU & $32\times32\times32$ \\
        Pool 1 & MaxPool $2\times2$; Dropout $p=0.3$ & $32\times16\times16$ \\
        Block 2 & Conv $32\rightarrow64$, $3\times3$; ReLU; Conv $64\rightarrow64$, $3\times3$; ReLU & $64\times16\times16$ \\
        Pool 2 & MaxPool $2\times2$; Dropout $p=0.3$ & $64\times8\times8$ \\
        Block 3 & Conv $64\rightarrow128$, $3\times3$; ReLU; MaxPool $2\times2$ & $128\times4\times4$ \\
        Classifier & Flatten; Linear $2048\rightarrow256$; ReLU; Dropout $p=0.5$; Linear $256\rightarrow10$ & 10 \\
        \bottomrule
    \end{tabularx}
\end{table}

\section{训练配置}

\begin{table}[H]
    \centering
    \caption{Baseline 训练设置}
    \begin{tabularx}{\textwidth}{lY}
        \toprule
        项目 & 设置 \\
        \midrule
        数据集 & CIFAR-10，本地路径 \texttt{data/} \\
        数据增强 & RandomCrop$(32,\mathrm{padding}=4)$，RandomHorizontalFlip \\
        归一化 & mean = $(0.4914, 0.4822, 0.4465)$，std = $(0.2470, 0.2435, 0.2616)$ \\
        损失函数 & CrossEntropyLoss \\
        优化器 & SGD，momentum = 0.9，Nesterov = True \\
        初始学习率 & 0.1 \\
        权重衰减 & $5\times10^{-4}$ \\
        学习率策略 & CosineAnnealingLR，$T_{\max}=80$ \\
        Batch size & 128 \\
        训练轮数 & 80 epochs \\
        随机种子 & 42 \\
        训练脚本 & \path{classic_CNN/train.py} \\
        模型文件 & \path{classic_CNN/model.py} \\
        \bottomrule
    \end{tabularx}
\end{table}

\section{实验结果}

完整训练 80 轮后，模型在测试集上取得的最佳分类准确率为 86.85\%，对应测试错误率为 13.15\%。
最优模型保存在 \path{classic_CNN/runs/baseline/best.pt}，训练日志保存在
\path{classic_CNN/runs/baseline/history.csv}。

\begin{table}[H]
    \centering
    \caption{ClassicCNN 在 CIFAR-10 上的结果}
    \begin{tabular}{ll}
        \toprule
        指标 & 数值 \\
        \midrule
        最佳 epoch & 80 \\
        Train loss & 0.4222 \\
        Train accuracy & 85.53\% \\
        Test loss & 0.3896 \\
        Best test accuracy & 86.85\% \\
        Test error & 13.15\% \\
        总训练时间 & 1796.67 s，约 29.94 min \\
        \bottomrule
    \end{tabular}
\end{table}

\begin{figure}[H]
    \centering
    \includegraphics[width=0.78\linewidth]{../runs/baseline/curves.png}
    \caption{ClassicCNN baseline 的训练与测试曲线}
\end{figure}

\end{document}
